% !TeX program = tectonic
\documentclass[11pt,a4paper]{article}
\usepackage{fontspec}
\usepackage[english]{babel}
\babelfont{rm}[Language=Default]{Liberation Serif}
\babelfont[Language=Default]{sf}{Carlito}
\babelfont[Language=Default]{tt}{DejaVu Sans Mono}
\usepackage{amsmath,amssymb,amsthm}
\usepackage{geometry}
\geometry{a4paper, top=2.4cm, bottom=2.4cm, left=2.4cm, right=2.4cm}
\usepackage{booktabs,tabularx,enumitem,xcolor,tcolorbox}
\tcbuselibrary{breakable,skins}
\definecolor{linkblue}{HTML}{1F4E79}
\definecolor{statgreen}{HTML}{2F6B3A}
\definecolor{goldacc}{HTML}{8B7E5A}
\usepackage[colorlinks=true,linkcolor=linkblue,urlcolor=linkblue,unicode=true]{hyperref}
\theoremstyle{plain}
\newtheorem{theorem}{Theorem}
\newtheorem{lemma}[theorem]{Lemma}
\newtheorem{proposition}[theorem]{Proposition}
\newtheorem{corollary}[theorem]{Corollary}
\theoremstyle{definition}
\newtheorem{definition}[theorem]{Definition}
\theoremstyle{remark}
\newtheorem{remark}[theorem]{Remark}
\newtcolorbox{statusbox}[1][]{colback=statgreen!5,colframe=statgreen!75,
  title={\textbf{Summary of results:} #1},fonttitle=\small,boxrule=0.7pt,arc=2pt,breakable}
\newtcolorbox{databox}[1][]{colback=goldacc!6,colframe=goldacc!80,
  title={\textbf{Protocol data:} #1},fonttitle=\small,boxrule=0.7pt,arc=2pt,breakable}
\newtcolorbox{verifybox}[1][]{colback=linkblue!5,colframe=linkblue!80,
  title={\textbf{Verification:} #1},fonttitle=\small,boxrule=0.7pt,arc=2pt,breakable}
\widowpenalty=10000
\clubpenalty=10000
\setlength{\parskip}{0.35em}
\newcommand{\CC}{\mathbb{C}}
\newcommand{\RR}{\mathbb{R}}
\newcommand{\ZZ}{\mathbb{Z}}
\newcommand{\QQ}{\mathbb{Q}}
\newcommand{\Ga}{\Gamma}
\newcommand{\Bfun}{\mathrm{B}}
\newcommand{\im}{\operatorname{Im}}
\newcommand{\re}{\operatorname{Re}}
\newcommand{\lcm}{\operatorname{lcm}}
\newcommand{\SNF}{\operatorname{SNF}}
\newcommand{\Jac}{\operatorname{Jac}}
\newcommand{\rank}{\operatorname{rank}}
\newcommand{\Ind}{\operatorname{Index}}
\newcommand{\disc}{\operatorname{disc}}
\begin{document}
\begin{center}
{\LARGE\bfseries Theorem 3: Polarization and Hodge Lattice Indices}\\[0.4em]
{\large the principal polarization and computability of block indices}\\[0.6em]
{\normalsize Isaev Iskhak Khamzatovich}\\[0.2em]
{\normalsize The <<Dynamic Principle>> Program --- the \texttt{hodge-laboratory}}\\[0.15em]
{\small Standalone theorem monograph · Монография 6/16}
\end{center}
\vspace{0.4em}
{\color{linkblue}\hrule height 0.8pt}
\vspace{0.8em}

\section{Setting}

The intersection form $J$ on $H_1(C_N,\ZZ)$ is not merely a
topological invariant: it is a polarization of the weight-1 Hodge
structure and defines the Hodge lattice indices per conductor
blocks --- integer invariants computable algorithmically.

\begin{definition}[polarization]
For closed forms $\widetilde E([\alpha],[\beta])
=\iint_{C_N}\alpha\wedge\beta$; the polarization form on
$H_1(C_N,\ZZ)$ is the intersection form $E:=J$.
\end{definition}

\begin{theorem}[$J$ is a weight-1 polarization]\label{th:pol}
With respect to the Hodge decomposition
$H^1=H^{1,0}\oplus H^{0,1}$ the Riemann--Hodge relations hold:
(1) $\widetilde E(\omega,\omega')=0$ for holomorphic forms;
(2) $-\frac1{2i}\widetilde E(\omega,\overline{\omega})
=\int_{C_N}|g|^2dxdy>0$. Hence $\Jac(C_N)$ with $J$ is a principally
polarized abelian variety.
\end{theorem}

\begin{proof}
(i) A curve has no nonzero holomorphic 2-forms, so
$\omega\wedge\omega'=0$ --- the first Riemann relation read on the
bilinear identity. (ii)
$\omega\wedge\overline{\omega}=-2i|g|^2dx\wedge dy$ gives
$-\frac1{2i}\widetilde E(\omega,\overline{\omega})
=\int_{C_N}|g|^2dxdy>0$. The unimodularity of $J$ gives the
principality.
\end{proof}

\begin{definition}[block indices]
For a conductor $d\mid N$: the block lattice $\Lambda_d$ --- the
isotypic component of characters with conductor $d$; the saturation
$\Lambda_d^{sat}$; the block index
$\Ind(\Lambda_d)=[\Lambda_d^{sat}:\Lambda_d]$. Block data: rank,
signature of $E|_{\Lambda_d}$, determinant, index.
\end{definition}

\begin{theorem}[computability of the indices]\label{th:indcomp}
All quantities are finite and computed algorithmically from the
closed forms: the saturation by Hermite--Smith reduction; the
signature and determinant from the combinatorial intersection
matrix filtered by characters; the index as a determinant ratio.
The counts per block: $1+6+84=91$ ($N=15$);
$1+6+9+30+84+276=406$ ($N=30$).
\end{theorem}

\begin{proof}
Finiteness: $H_1(C_N,\ZZ)$ is free abelian of rank $2g$.
Computability: the isotypic component is defined over
$\QQ(\zeta_N)$; its intersection with the integer lattice is an
integer sublattice; the saturation basis by the Hermite--Smith
algorithm on the $\mu_N\times\mu_N$-action table (sheet
combinatorics). The restricted polarization: $E=J$, a combinatorial
matrix. The census sums match the genera (V1).
\end{proof}

\begin{databox}[numbers --- certificate H]
The stand polarization type (SNF): $(1,1,5,5,15,15,15,15)$; the
product of invariant factors $=1265625=3^4\cdot5^6=1125^2$ ---
exact integer agreement with the discriminant; $\mathrm{vol}_h=1125$.
\end{databox}

\begin{statusbox}[summary]
The polarization is proved (Riemann--Hodge); the computability is
proved algorithmically; certificate H implements the computation
with exact integer certification --- the indices are not only
computable but computed.
\end{statusbox}

\begin{verifybox}[reproduction]
\texttt{python3 laboratory.py} $\to$ ``Certificates $\to$ H1'';
Lean: \texttt{snf\_product}, \texttt{stand\_disc}.
\end{verifybox}

\vfill
{\color{linkblue}\hrule height 0.6pt}
\vspace{0.5em}
{\small\color{gray}
License: individual exclusive license of Isaev Iskhak Khamzatovich.
All rights to the computations, formulas, and texts belong to the author.
Verification: \texttt{python3 laboratory.py} (``Stands''/``Certificates''),
Lean: \texttt{verification/lean/HodgeLaboratory.lean}.
}
\end{document}
